\documentclass[a4paper,12pt]{scrartcl}
\usepackage[utf8x]{inputenc}
\usepackage[T1]{fontenc} % avec T1 comme option  d'encodage c'est ben mieux, surtout pour taper du français.
%\usepackage{lmodern,textcomp} % fortement conseillé pour les pdf. On peut mettre autre chose : kpfonts, fourier,...
\usepackage[french]{babel, varioref} %Sans ça les guillemets, amarchpo
\usepackage{amsmath}
\usepackage{multicol}
\usepackage{amssymb}
\usepackage{tkz-tab}
\usepackage{exercice_sheet}
\flushbottom
\raggedbottom

%\trait
%\section*{}
%\exo{}
%\question{}
%\subquestion{}

\date{}

% Title Page
\title{Équations du second degré}

\author{}

\begin{document}

\selectlanguage{french}

\maketitle

\tableofcontents

\section{Définition}

Une équation du second degré à une inconnue $x$ est une équation pouvant se ramener à la forme standard suivante: \answer{$ax^2 + bx + c = 0$} où $a$ est un réel non nul, $b$ et $c$ sont des réels quelconques.

Le qualificatif \emph{second degré} vient du fait que l'inconnue est portée au maximum à la puissance $2$ dans l'équation.

\exemple{}

\begin{itemize}
\item L'équation $3x^2 - 2x + 4 = 0$ est une équation du 2\textsuperscript{nd} degré où $a = \ldots$, $b = \ldots$ et $c = \ldots$ 
\item L'équation $x^2 - x + 1 = 0$ est une équation du 2\textsuperscript{nd} degré où $a = \ldots$, $b = \ldots$ et $c = \ldots$ 
\item L'équation $5x^2 +3x = 0$ est une équation du 2\textsuperscript{nd} degré où $a = \ldots$, $b = \ldots$ et $c = \ldots$ 
\item L'équation $1.5x^2 + 3/2 = 0$ est une équation du 2\textsuperscript{nd} degré où $a = \ldots$, $b = \ldots$ et $c = \ldots$ 
\end{itemize}

\exemple{}
$5x^2 + 3x = 4x^2 + 8x - 6$ est une équation du second degré car elle peut se
ramener à la forme standard en appliquant les règles générales régissant les équations:

$5x^2 + 3x = 4x^2 + 8x - 6$

$\Leftrightarrow 5x^2 - 4x^2 + 3x - 8x +6 = 0$

$\Leftrightarrow x^2 + - 5x +6 = 0$

\section{Méthode algébrique de résolution}\label{methode}

On peut montrer qu’une équation du second degré ne peut avoir que zéro, une ou deux solutions. 

\begin{itemize}
\item on ramène l’équation sous sa forme standard: $ax^2 + bx + c = 0$;
\item on identifie les valeurs des nombres $a$, $b$ et $c$;
\item on calcule le nombre $\Delta = b^2 - 4ac$. Ce nombre est appelé le \emph{discriminant}.
\item on peut conclure suivant la valeur de $\Delta$:

\begin{itemize}
\item $\Delta < 0$, l'équation n'a pas de solution: $S = \emptyset$.
\item $\Delta = 0$, a une unique solution: $x_1 = \dfrac{-b}{2a}$.

On écrit l'ensemble des solutions:

$S = \left\lbrace \dfrac{-b}{2a} \right\rbrace$.

\item $\Delta > 0$, l'équation a deux solutions: $x_1 = \dfrac{-b - \sqrt{\Delta}}{2a}$ et $x_2 = \dfrac{-b + \sqrt{\Delta}}{2a}$.

On écrit l'ensemble des solutions:

$S = \left\lbrace \dfrac{-b - \sqrt{\Delta}}{2a} ; \dfrac{-b + \sqrt{\Delta}}{2a} \right\rbrace$.

\end{itemize}
\end{itemize}

À noter que l'on appelle ces solutions les \textbf{racines} du polynôme.

%\subsection*{En résumé}

%La méthode de résolution se synthétise comme suit:

%Mettre algorithme (http://www.texample.net/tikz/examples/simple-flow-chart/)

\section{Quelques exemples}

\exemple{}
$x^2 - 5x + 6 = 0$

\begin{itemize}
\item Ici, l'équation est déjà sous sa forme standard. $a = \ldots$, $b = \ldots$ et $c = \ldots$
\item on calcule $\Delta = b^2 - 4ac$, soit $\Delta = \ldots\ldots\ldots = \ldots$
\item $\Delta > 0$, il y a donc 2 solutions: 

$x_1 = \dfrac{-b - \sqrt{\Delta}}{2a}$ et $x_2 = \dfrac{-b + \sqrt{\Delta}}{2a}$.

Soit :

$x_1 = \dfrac{-\ldots - \sqrt{\ldots}}{2\times \ldots}$ et $x_2 = \dfrac{-\ldots + \sqrt{\ldots}}{2\times \ldots}$.

D'où $x_1 = \ldots$ et $x_2 = \ldots$

\item On peut conclure: $S = \left\lbrace 2;3 \right\rbrace$.

\end{itemize}


\exemple{}
$5^2 + 2x + 4 = 2x^2 - 1$

\begin{itemize}
\item On ramène l'équation sous sa forme standard:

$5^2 + 2x + 4 = 2x^2 - 1$

$\Leftrightarrow 5^2 + 2x + 4 -  2x^2 + 1= 0$

$\Leftrightarrow 3^2 + 2x + 5 = 0$

\item On reconnaît: $a = \ldots$, $b = \ldots$ et $c = \ldots$
\item on calcule $\Delta = b^2 - 4ac$, soit $\Delta = \ldots\ldots\ldots = \ldots$
\item $\Delta < 0$, il n'y a donc pas de solution: 

$S = \emptyset$.

\end{itemize}

\exemple{}
$x^2 +2x + 1 = 0$

\begin{itemize}
\item L'équation est déjà sous sa forme standard. $a = \ldots$, $b = \ldots$ et $c = \ldots$
\item on calcule $\Delta = b^2 - 4ac$, soit $\Delta = \ldots\ldots\ldots = \ldots$
\item $\Delta = 0$, il y a donc 1 solution: 

$x_1 = \dfrac{-b }{2a}$.

Soit :

$x_1 = \dfrac{-\ldots}{2\times \ldots}$

D'où $x_1 = \ldots$.

\item On peut conclure: $S = \left\lbrace \ldots \right\rbrace$

\end{itemize}

\section{Cas particuliers}

La méthode décrite précédemment est générale et peut s'appliquer à \emph{toutes} les équations du second degré.

Il existe néanmoins des cas particuliers pour lesquels on peut résoudre plus simplement.

\subsection{Cas où $b = 0$, $ax^2 + c = 0$}

\exemple{}
$2x^2 + 4 = 0$

$\Leftrightarrow 2x^2 = -4$

$\Leftrightarrow x^2 = -2$

Un carré ne peut être négatif, il s'agit donc d'une équation sans solution: $S = \emptyset$.

\exemple{}
$2x^2 - 8 = 0$

$\Leftrightarrow 2x^2 = 8$

$\Leftrightarrow x^2 = 4$

$\Leftrightarrow 
\left\lbrace
\begin{array}{ll}
   x = -2 \\
   \mbox{ou} \\
   x = 2 \\
\end{array}
\right.
$ 

$S = \left\lbrace -2 ; 2 \right\rbrace$

\subsection{Cas où $c = 0$, $ax^2 + bx = 0$}

Cette équation se factorise par $x$ : 

\exemple{}
$5x^2 - 6x = 0$

$\Leftrightarrow x(5x-6) = 0$

Or, on sait que lorsque le produit de plusieurs facteurs est égal à $0$, alors au moins un de ces facteurs est égal à $0$. On obtient donc :

$\left\lbrace
\begin{array}{ll}
   x = 0 \\
   \mbox{ou} \\
   5x-6 = 0 \\
\end{array}
\right.
$ 

$
\Leftrightarrow
\left\lbrace
\begin{array}{ll}
   x = 0 \\
   \mbox{ou} \\
   x = \dfrac{6}{5} \\
\end{array}
\right.
$ 

L'ensemble des solutions est donc $S = \left\lbrace 0 ; \dfrac{6}{5} \right\rbrace$

\subsection{Cas de l'équation produit $(a_1 x + b_1)(a_2 x + b_2) = 0$}

\exemple{}
$(2x +1)(3x -5) = 0$ 

De la même manière que dans l'exemple précédent, cette équation est équivalente à:

$\left\lbrace
\begin{array}{ll}
   2x +1 = 0 \\
   \mbox{ou} \\
   3x -5 = 0 \\
\end{array}
\right.
$ 

$
\Leftrightarrow
\left\lbrace
\begin{array}{ll}
   x = -\dfrac{1}{2} \\
   \mbox{ou} \\
   x = \dfrac{5}{3} \\
\end{array}
\right.
$ 

L'ensemble des solutions est donc $S = \left\lbrace -\dfrac{1}{2} ; \dfrac{5}{3} \right\rbrace$

\section{Factorisation}

\subsection{Méthode}

Une fonction polynomiale du second degré $P(x)$ est une expression qui, sous sa forme développée\footnote{C'est-à-dire la forme où les dernières opérations réalisées sont des additions et/ou des soustractions.} s'écrit:

\Answer{$P(x) = ax^2 + bx + c$}

On recherche les solutions de l'équation $P(x) = 0$ comme vu au paragraphe \vref{methode}. On admet qu'on peut factoriser le polynôme à l'aide des solutions de cette équation (voir paragraphe \vref{methode}) comme suit: 

\begin{itemize}
\item $\Delta > 0$, l'équation $P(x) = 0$ a deux solutions $x_1$ et $x_2$ et on peut factoriser le polynôme ainsi: \answer{$P(x) = a(x-x_1)(x-x_2)$}

\item $\Delta = 0$, l'équation $P(x) = 0$ a une solution $x_1$ et on peut factoriser le polynôme ainsi: \answer{$P(x) = a(x-x_1)^2$}

\item $\Delta < 0$, l'équation $P(x) = 0$ n'a pas de solution dans $\mathbb{R}$ et on ne peut factoriser le polynôme
\end{itemize}

\subsection{Exemples}

\exemple{}
Factoriser le polynôme $P(x) = x^2 - 5x + 6$

On calcule $\Delta$. Ici, $\Delta = 1 > 0$. L'équation a donc 2 solutions qui sont $x_1 = 2$ et $x_2 = 3$. De plus, $a = 1$. Le polynôme se factorise donc\footnote{Pour se convaincre que la méthode fonctionne, il suffit de prendre ce résultat et de le développer, on tombe sur le polynôme donné au départ...} de la façon suivante: $P(x) = a(x - x_1) (x - x_2)$

On obtient donc $P(x) = \ldots\ldots\ldots\ldots$

\exemple{}
Factoriser le polynôme $P(x) = 3x^2 + 2x + 5$.

$\Delta = -56 < 0$ l'équation $P(x) = 0$ n'a donc pas de solution. 

Le polynôme n'est pas factorisable dans $\mathbb{R}$.

\exemple{}
Factoriser le polynôme $P(x) = x^2 - 2x + 1$.

$\Delta = 0$. L'équation $P(x) = 0$ a donc une unique solution $x_1 = 1$. 

Le polynôme se factorise donc de la façon suivante: $P(x) = a(x - x_1)^2$. 

De plus, $a = 1$.

On obtient donc $P(x) = \ldots\ldots\ldots\ldots$

\section*{Exercices}

\exo[1]{Résoudre les équations suivantes:}

\question{}
$x^2 - 2x -15 = 0$

\question{}
$4x (x +1) = -1$

\question{}
$(2x -1)(x + 1) = -6$

\question{}
$2x^2 + x = 0$

\question{}
$(x +2)(x -1) = 0$

\exo[2]{}
La \vref{tranchee} représente une tranchée rectangulaire ABCD d'une
profondeur maximale de 10 m.

\begin{figure}
\begin{center}
\includegraphics[width=.5\textwidth]{pics/tranchee.pdf}
\end{center}
\caption{Schéma de la tranchée}
\label{tranchee}
\end{figure}

\question{}
Déterminer l’aire du rectangle ABCD pour $x = 6 m$.

\question{}
Exprimer l’aire $S(x)$ du rectangle ABCD en fonction de $x$.

On considère dans la suite de l’exercice la fonction $S$ définie sur l’intervalle $[0 ;10]$ par :
$S(x) = -2x^2 + 20x$.

\question{}
Compléter le tableau de valeurs suivant:

\begin{table}[h]
\begin{center}
\begin{tabular}{|c|c|c|c|c|c|c|c|c|c|c|c|}
\hline
$x$    & 0 & 1  & 2  & 3  & 4  & 5  & 6  & 7  & 8  & 9  & 10 \\ \hline
$S(x)$ & \hspace{0.6cm} & \hspace{0.6cm} & \hspace{0.6cm} & \hspace{0.6cm} & \hspace{0.6cm} & \hspace{0.6cm} & \hspace{0.6cm} & \hspace{0.6cm} & \hspace{0.6cm} & \hspace{0.6cm} & \hspace{0.6cm} \\ \hline
\end{tabular}
\end{center}
\end{table}

\question{}
Tracer la courbe $\mathcal{C}_S$, représentation graphique de la fonction $S$ dans le repère fourni.

Échelle :

\begin{itemize}
\item axe des abscisses : 1 cm pour 1 m.
\item axe des ordonnées : 2 cm pour 10 m² .
\end{itemize}

\question{}
Variations de la fonction.

\subquestion{}
Établir le tableau de variation de la fonction $S$.

\subquestion{}
Déterminer la valeur de $x$ pour laquelle $S(x)$ est maximale. Quelle est cette valeur maximale ?

\question{}
Résolution d'équation

\subquestion{}
Déterminer graphiquement les valeurs de $x$ pour lesquelles l’aire du rectangle
ABCD est égale à 37.5 m². Les traits de construction devront figurer sur le
graphique.

\subquestion{}
Retrouver les résultats précédents en résolvant l’équation : $-2x^2 + 20x = 37.5$.

%\exo{Résoudre les équations suivantes :}
%
%\question{$x^2 - 2x - 15 = 0$}
%
%
%\question{$2x^2 + x + 5 = 0$}
%
%
%\question{$4x^2 + 4x + 1 = 0$}
%
%\question{$8x^2 + 2x = 15 + 3x + 2x^2$}
%
%
%\question{$(x + 2) (x - 1) = 0$}
%
%
%\question{$x^2 + 2x = 0$}





\end{document}
